\documentclass[twocolumn]{iopjournal}

\usepackage{amsmath,amssymb,amsfonts}

\fancyhfoffset[L]{0mm}
\setlength{\marginparwidth}{0pt}
\setlength{\marginparsep}{0pt}

% ==========================================
% ZENODO PREPRINT HEADERS (ACTIVE)
% ==========================================
\renewcommand{\articletype}[1]{%
  {\vspace*{-8mm}\noindent \Large \sf Preprint}\par
  \vspace*{6mm}\noindent{\scriptsize \sf{\bfseries \MakeUppercase{#1}}}}

\fancyhead[L]{}
\fancyhead[R]{}

\begin{document}

\articletype{Supplementary Material}

\title{Supplementary Material: From the Lorentz Signature\\to the Tetrahedral Frame}

\author{Stephen Nelson}

\affil{Field Dynamics Research}

\begin{abstract}
In a Lorentzian local frame with signature $(-,+,+,+)$, the projective null cone is naturally identified with the unit sphere $S^2$ of spatial directions. Requiring that a discrete four-direction realisation introduce no net preferred spatial direction and treat all selected directions equivalently uniquely fixes the regular tetrahedral frame. The result contains no free parameter and supplies the four directions used in the main letter's half-space construction.
\end{abstract}

% ----------------------------------------------------------------------

\section*{Derivation}

\subsection*{Step 0: The tetrad exists}

Any Lorentzian 4-manifold $(M, g)$ admits a local orthonormal tetrad $e_{(a)}$ satisfying
\begin{equation}
g(e_{(a)}, e_{(b)}) = \eta_{ab} = \mathrm{diag}(-1,+1,+1,+1).
\end{equation}
This is not a choice. It is a theorem of differential geometry. The tetrad has four legs: one timelike, three spacelike.

\subsection*{Step 1: The null cone maps directions to a sphere}

Every future-directed null direction in the local frame is represented, up to positive rescaling, by
\begin{equation}
k \sim e_{(0)} + u^i\, e_{(i)}, \qquad u \in \mathbb{R}^3.
\end{equation}
The null condition $g(k,k) = 0$ gives
\begin{equation}
-1 + |u|^2 = 0,
\end{equation}
so $|u| = 1$. The projective null cone of $(-,+,+,+)$ is in bijection with the unit sphere $S^2$.

\subsection*{Step 2: Isotropy fixes the admissible four-direction set}

The three spatial entries of $\eta_{ab} = \mathrm{diag}(-1,+1,+1,+1)$ are all $+1$. The spatial sector is therefore Euclidean, with isotropy group $SO(3)$; no spatial axis is metrically privileged. A discrete set of directions on $S^2$ that weighted one axis more than another would introduce a distinction the metric does not contain.

% NOTE: The number four is not inherited from the tetrad's four legs.
% It comes from the spatial dimension: in D=3, the minimal closed
% solid (the simplex) has D+1=4 faces and hence four outward normals.
% The tetrad also has four legs (one time, three space) because
% spacetime is D+1=3+1 dimensional; both fours trace back to D=3,
% but through different paths.
In $D = 3$ spatial dimensions, the minimal number of bounding planes that can enclose a finite region is $D + 1 = 4$. Any discrete direction set that faithfully represents the spatial sector must therefore contain at least four directions $u_I \in S^2$ ($I = 1,\dots,4$). Because the Lorentz signature is exactly balanced in its spatial entries, this four-direction set must not introduce anisotropy absent from the parent frame; its discrete multipole moments must respect the continuous $SO(3)$ of the spatial sector:
\begin{itemize}
\item Unit norm: $|u_I| = 1$ for all $I$ \;\;(each direction lies on the null sphere).
\item Vanishing first moment: $\displaystyle\sum_{I=1}^{4} u_I = 0$ \;\;(no net directional bias).
\item Isotropic second moment: $\displaystyle T^{ij} := \sum_{I=1}^{4} u_I^i\, u_I^j \;\propto\; \delta^{ij}$ \;\;(no preferred axis or plane).
\end{itemize}
These are the discrete translations of $SO(3)$ invariance at the first two multipole orders, and for a four-point set they are exhaustive: four unit vectors with fixed centroid and Gram matrix are determined up to $O(3)$, so the three conditions fully constrain the configuration. To see this, note that the Gram matrix $G_{IJ} = u_I \cdot u_J$ is a $4 \times 4$ symmetric matrix whose diagonal entries are $G_{II} = 1$ (unit norm), whose kernel contains $\mathbf{1}$ (zero centroid), and whose nonzero eigenvalues are all equal (isotropic second moment). These constraints force $G_{IJ} = -1/3$ for all $I \neq J$, leaving a single common pairwise inner product $c = -1/3$. The problem therefore reduces to classifying four-point subsets of $S^2$ satisfying these conditions.

\subsection*{Proposition}

Under the three conditions above, any four-point subset $\{u_I\}_{I=1}^4 \subset S^2$ is, up to $O(3)$, the regular tetrahedral frame.

\subsection*{Proof}

\noindent\textit{Step 3: Solve.}\quad Square the zero-centroid condition:
\begin{equation}
0 = \left|\sum_{I=1}^{4} u_I\right|^2 = \sum_{I=1}^{4} |u_I|^2 + 2\!\!\sum_{1 \leq I < J \leq 4} u_I \cdot u_J.
\end{equation}
Since $|u_I|^2 = 1$ for all $I$, and there are $\binom{4}{2} = 6$ pairs,
\begin{equation}
0 = 4 + 12c \quad \implies \quad c = -\frac{1}{3}.
\end{equation}
There is no free parameter. Once the four-direction set is required to introduce no net preferred spatial direction and to treat all members equivalently, the inner product is fixed to $c = -1/3$.

\noindent\textit{Step 4: Identify.}\quad Four unit vectors in $\mathbb{R}^3$ with constant pairwise inner product $-1/3$ are the vertices of a regular tetrahedron inscribed in $S^2$, unique up to $O(3)$~\cite{Coxeter1973}. The explicit realisation is:
\begin{align}
u_1 &= \tfrac{1}{\sqrt{3}}\,(+1,+1,+1), \nonumber \\
u_2 &= \tfrac{1}{\sqrt{3}}\,(+1,-1,-1), \nonumber \\
u_3 &= \tfrac{1}{\sqrt{3}}\,(-1,+1,-1), \nonumber \\
u_4 &= \tfrac{1}{\sqrt{3}}\,(-1,-1,+1).
\end{align}
These satisfy $|u_I| = 1$, $\sum_I u_I = 0$, and $u_I \cdot u_J = -1/3$ for $I \neq J$. They are the four projections $v_1$ through $v_4$ in the main letter, normalised. The four directions are not freely chosen; they are uniquely determined up to $O(3)$. \hfill$\square$

\subsection*{Step 5: Close}

The four half-space inequalities in the main letter are defined by the planes normal to these four directions. The chamber decomposition proved in the main letter therefore inherits its structure from a tetrahedral frame whose isotropy is not an independent condition but the discrete expression of the Lorentz signature's exactly balanced spatial sector.

% ----------------------------------------------------------------------

\section*{Summary}

The $(-,+,+,+)$ signature maps null directions to $S^2$. Requiring that a discrete four-direction set on $S^2$ introduce no net preferred spatial direction uniquely fixes the regular tetrahedral frame with $c = -1/3$. Its four face normals are the four projections used in the main letter. The boundedness proof does the rest.

% ----------------------------------------------------------------------

\begin{thebibliography}{99}

\bibitem{Coxeter1973}
H.~S.~M.~Coxeter,
\textit{Regular Polytopes},
3rd ed.\ (Dover, New York, 1973).

\end{thebibliography}

\end{document}
